\subsection{辛群。}

假设环 A 交换。若 \(\Phi\) 是 E 上的一个交错双线性型，则使 \(\Phi\) 不变的模 E 的自同构称为关于 \(\Phi\) 的辛自同构（或辛变换），它们构成一个群，称为与 \(\Phi\) 相伴的辛群；它有时记作 \(\mathbf{Sp}(\Phi)\) 。

特别地，考虑模 \(E = A^{2m}\) 上的交错双线性型 \(\Phi_{0}\) ，它关于 E 的典范基 \((e_{i})\) 的矩阵是

\[
R _ {m} = \left( \begin{array}{c c} 0 & I _ {m} \\ - I _ {m} & 0 \end{array} \right).
\]

关于 \(\Phi_{0}\) 的辛自同构与辛群分别简称为（A 上的）\(2m\) 变量的辛自同构与辛群；这个群记作 \(\mathbf{Sp}(2m, A)\) 或 \(\mathbf{Sp}_{2m}(A)\) 。把辛自同构的每个矩阵 \(A\) （关于典范基 \((e)_i\) ）称为辛矩阵。这样的矩阵是可逆的，并且由 § 1，第 10 小节的公式 (48)，满足关系

\[
{ } ^ { t } A . R _ { m } . A = R _ { m } .\tag{9}
\]

反之，若 A 上的 2m 阶方阵 A 满足 (9)，则它是辛的：实际上只需证明它可逆；而 (9) 蕴含 Pf(R \(_{m}\) ) = Pf( \(^{t}\) A . R \(_{m}\) . A) = (det A)Pf(R \(_{m}\) )（第 2 小节命题 1），因而 det A = 1。我们同时证明了下面的命题：

命题 3. --- 辛矩阵的行列式等于 1。

若 A 是一个域，\(\Phi\) 是 A 上偶数维 \(2m\) 向量空间 E 上的非退化交错双线性型，则由定理 1 的推论 1，与 \(\Phi\) 相伴的辛群同构于 \(\mathbf{Sp}(2m, \mathrm{A})\) 。

习题. --- ¶ 1) 设 A 是一个主交换环，E 是有限维数为 n 的自由 A-模，Φ 是 E 上的交错双线性型；第 1 小节定理 1 中定义的理想 Aαi（1 ≤ i ≤ r）称为 Φ 的不变因子。

\begin{enumerate}
\def\labelenumi{\alph{enumi})}
\item
  设 F 是 E 的一个子模，\(\Phi_{\mathrm{F}}\) 是 \(\Phi\) 在 F × F 上的限制。证明：若 A \(\beta_{i}\) （ \(1 \leqslant i \leqslant s\) ）是 \(\Phi_{\mathrm{F}}\) 的不变因子（其中 \(\beta_{i}\) 整除 \(\beta_{i+1}\) ），则 \(s \leqslant r\) ，且 \(\beta_{i}\) 是 \(\alpha_{i}\) 的倍式（对 \(1 \leqslant i \leqslant s\) ）。（化归到 \(r = s = n/2\) 的情形，并利用第 VII 章，§ 4 的习题 9 b) 与 9 c)。）
\item
  设 \(E_{1}\) 是另一个有限维自由 A-模，\(\Phi_{1}\) 是其上的一个交错双线性型，模与不变因子为 \(E_{1}, A\gamma_{1}, \ldots, A\gamma_{s}\) （ \(\gamma_{i}\) 整除 \(\gamma_{i+1}\) ）。要使 \(\Phi_{1}\) 是 \(\Phi\) 在 \(E_{1}\) 到 E 中的某个线性映射下的原像，必须且只需 \(s \leqslant r\) ，且 \(\gamma_{i}\) 是 \(\alpha_{i}\) 的倍式（对 \(1 \leqslant i \leqslant s\) ）。（利用 a) 与第 VII 章，§ 4，第 5 小节的命题 4）。
\item
  设 F、G 是 E 的两个子模，使得 \(F^{0}\) （相应地，\(G^{0}\) ）是 F（相应地，G）在 E 中的补。若 \(\Phi\) 在 F 与 G 上的限制等价，证明 \(\Phi\) 在 \(F^{0}\) 与 \(G^{0}\) 上的限制也是如此，并且存在 E 的自同构，它保持 \(\Phi\) 不变且把 F 映到 G。
\item
  举出 E 的两个维数为 2 的子模 F、G 的例子，使得 F 与 G 在 E 中都有补，且限制 \(\Phi_{\mathrm{F}}\) 与 \(\Phi_{\mathrm{G}}\) 等价，但不存在 E 的自同构保持 \(\Phi\) 不变且把 F 映到 G（取 \(n = 4\) ）。
\end{enumerate}

\begin{enumerate}
\def\labelenumi{\arabic{enumi})}
\setcounter{enumi}{1}
\tightlist
\item
  设 \(\Phi\) 是有限维向量空间 E 上的一个交错双线性型。证明：对 E 的每个子空间 M，差 dim M - dim (M ∩ M⁰) 是偶数。（先考虑 Φ 非退化的情形。）
\end{enumerate}

¶ 3) 设 E 是域 A 上维数 n = 2m 为偶数的向量空间；设 Φ 与 Ψ 是 E 上的两个交错双线性型；假设 Ψ 非退化。设 u 与 v 分别是在右边与 Φ 及 Ψ 相伴的从 E 到 E* 的线性映射；v 是从 E 到 E* 上的同构；令 w = v⁻¹ ∘ u，于是 w 是 E 的自同态。

\begin{enumerate}
\def\labelenumi{\alph{enumi})}
\item
  令 \(\mathbf{M}_{\mathbf{0}} = \mathbf{E}\) ，并递归地令 \(\mathbf{M}_{k+1} = w(\mathbf{M}_{k})\) （对 \(k > 0\) ）。证明：若 \(\mathbf{M}_{k}^{\prime}\) 是 \(\mathbf{M}_{k}\) 关于 \(\Phi\) 的正交子空间，则 \(\mathbf{M}_{k+1}\) 是 \(\mathbf{M}_{k}^{\prime}\) 关于 \(\Psi\) 的正交子空间。
\item
  设 \(n_{0}\) 是 \(\mathbf{M}_{0}^{\prime}\) 的维数；令 \(m_{0}=0\) ，且对 \(k\geqslant1\) ，以 \(m_{k}\) 表示 \(\mathbf{M}_{k}\cap\mathbf{M}_{0}^{\prime}\) 的维数。证明：对 \(k\geqslant1\) ，\(\mathbf{M}_{k}\) 的维数是 \(n-n_{0}-(m_{1}+\ldots+m_{k-1})\) ，而 \(\mathbf{M}_{k}^{\prime}\) 的维数是 \(n_{0}+(m_{1}+\ldots+m_{k})\) 。
\item
  证明：对一切 \(k \geqslant 0\) ，\(\mathbf{M}_k \cap \mathbf{M}_k'\) 的维数是 \(m_k + m_{k+1} + \ldots + m_{2k}\) ，而 \(\mathbf{M}_k' \cap \mathbf{M}_{k+1}\) 的维数是 \(m_{k+1} + m_{k+2} + \ldots + m_{2k+1}\) （应用 § 3 的习题 2 b) 来计算 \(\mathbf{M}_h \cap \mathbf{M}_{2k-h}'\) 与 \(\mathbf{M}_h' \cap \mathbf{M}_{2k+1-h}\) 的维数，对 \(h\) 作归纳）。
\item
  由 c) 推出数 \(m_{k}\) 都是偶数（利用习题 2）。
\item
  由 d) 得出：自同态 w 的对应于特征根 \(\lambda = 0\) 且有给定次数的初等因子之个数是偶数（参见第 VII 章，§ 5，习题 20）。
\end{enumerate}

\begin{enumerate}
\def\labelenumi{\arabic{enumi})}
\setcounter{enumi}{3}
\item
  设 E 是域 A 上的向量空间，具有可数基 \((e_n)_{n\geqslant 1}\) ，\(\Phi\) 是 E 上的非退化交错型。证明：E 中存在一个基 \((a_n)\) ，使得 \(\Phi(a_{2n-1}, a_{2n}) = 1\) （对一切 \(n\geqslant 1\) ），且 \(\Phi(a_i, a_j) = 0\) （对任何其他满足 \(i < j\) 的指标对）（按 § 4 的习题 13 的方式推理）。
\item
  对任意交错矩阵 \(X = (x_{ij})\) （偶数阶 \(n = 2m\) ，在交换环上）以及每个指标 \(i\) ，证明有 \(\operatorname{Pf}(X) = \sum_{j=1}^{n} (-1)^{i+j-1} \operatorname{Pf}(X_{ij}) x_{ij}\) ，其中 \(X_{ij}\) 是 \(n - 2\) 阶矩阵，它由从 \(X\) 删去指标为 \(i\) 与 \(j\) 的行和列而得到。
\item
  设 M 是交换环上的 m 阶方阵。证明：若令
\end{enumerate}

\[
R = \left( \begin{array}{c c} 0 & M \\ - ^ {t} M & 0 \end{array} \right)
\]

则有 \(\operatorname{Pf}(R) = \det M\) （先在 \(M\) 可逆时证明它，利用公式 (7)）。

\begin{enumerate}
\def\labelenumi{\arabic{enumi})}
\setcounter{enumi}{6}
\item
  设 A 是一个交换域，\(\mathfrak{A}(A)\) 是 A 上 \(2m\) 阶交错矩阵所成的集；设 I 是从 \(\mathfrak{A}(A)\) 到 A 中的一个映射，使得对每个矩阵 \(R\in\mathfrak{A}(A)\) 以及 A 上每个矩阵 \(P\) （阶为 \(2m\) ），有 \(\mathrm{I}(^{t}PRP)=(\det P)^{h}\mathrm{I}(R)\) ，其中 \(h\) 是一个有理整数。证明 \(\mathrm{I}(R)=c(\mathrm{Pf}R)^{h}\) ，其中 \(c\in\mathrm{A}\) （利用公式 (7) 与第 1 小节的定理 1）。
\item
  设 \(P\) 、\(Q\) 是交换环 A 上两个偶数阶 \(2m\) 的交错方阵。令 \(\varphi(X) = \mathrm{Pf}(P - XQ)\) ；证明：若 \(Q\) 可逆，则有 \(\varphi(Q^{-1}P)=0\) 。（先考虑 A 是域的情形；然后由习题 3 推出矩阵 \(Q^{-1}P\) 的极小多项式整除 \(\varphi(X)\) ，办法是转到 A 的一个代数闭扩张，并注意 \(\varphi^{2}\) 与 \(Q^{-1}P\) 的特征多项式只差一个数量因子。）
\end{enumerate}

¶ 9) 设 E 是交换域 A 上维数 n 的向量空间，Φ 是 E 上的一个交错双线性型，Ψ 是 E 上的一个埃尔米特半双线性型（关于 A 的一个对合自同构），P 与 Q 分别是 Φ 与 Ψ 关于 E 的同一个基的矩阵。

\begin{enumerate}
\def\labelenumi{\alph{enumi})}
\item
  我们假设 \(\Psi\) 非退化。证明：\(Q^{-1}P\) 的对应于特征根 0 且具有相同偶次数的初等因子之个数是偶数（习题 3 的方法）。
\item
  我们假设 \(\Phi\) 非退化（这蕴含 \(n\) 是偶数）。证明：\(P^{-1}Q\) 的对应于特征根 0 且具有相同奇次数的初等因子之个数是偶数（同样的方法）。
\end{enumerate}

\begin{enumerate}
\def\labelenumi{\arabic{enumi})}
\setcounter{enumi}{9}
\tightlist
\item
  以 \(\omega_{2m}(\mathbf{F}_q)\) 表示辛群 \(\mathbf{Sp}(2m, \mathbf{F}_q)\) （在有限域 \(\mathbf{F}_q\) 上）的阶。证明：若 \(h_{2m}\) 是向量对 \((x, y)\) （属于 \(\mathbf{F}_q^{2m}\) ）中使 \(\Phi_0(x, y) = 1\) 的个数（第 3 小节的记号），则有 \(\omega_{2m}(\mathbf{F}_q) = h_{2m}\omega_{2m-2}(\mathbf{F}_q)\) ；由此推出
\end{enumerate}

\[
\omega_ {2 m} (\mathbf {F} _ {q}) = (q ^ {2 m} - 1) q ^ {2 m - 1} (q ^ {2 m - 2} - 1) q ^ {2 m - 3} \dots (q ^ {2} - 1) q.
\]

\begin{enumerate}
\def\labelenumi{\arabic{enumi})}
\setcounter{enumi}{10}
\tightlist
\item
  设 A 是一个交换域。证明：每个变换 \(u\) （属于辛群 \(\mathbf{Sp}(2m, \mathrm{A})\) ）都是属于这个群的一些平延之积（称为辛平延；参见 § 4，习题 6）。（对 \(m\) 作归纳，办法是证明：若 \(x, y\) 是 \(\mathrm{E} = \mathrm{A}^{2m}\) 的两个不正交的向量，则存在辛平延的一个乘积 \(\nu\) ，使得 \(\nu u\) 保持 \(x\) 与 \(y\) 不变）。由此给出第 3 小节命题 3 的一个新证明。
\end{enumerate}

*12) 设 A 是一个交换域，E 是 A 上偶数维 \(n=2m\) 的向量空间，\(\Phi\) 是 E 上的一个非退化交错双线性型。证明：对每个相似变换 \(u\) （关于型 \(\Phi\) ； \(\S 6\) ，第 5 小节），若其乘子为 \(\alpha\) ，我们有 det \(u=\alpha^{m}\) （利用公式 (7)）。*

¶ 13) 我们假设 A 是特征 0 的交换域，E 是 A 上维数 2m 的向量空间，Φ 是 E 上的一个非退化交错双线性型。我们把 Φ 的逆形式 \(\widehat{\Phi}\) 等同于一个双向量 \(\Gamma\in\bigwedge E\) ，使得对 E 的任何辛基 \((e_{i})_{1\leqslant i\leqslant2m}\) （关于 \(\Phi\) 的），其编号适合 \(\Phi(e_{i},e_{j})=\Phi(e_{m+i},e_{m+j})=0\) ，\(\Phi(e_{i},e_{m+j})=\delta_{ij}(1\leqslant i\leqslant m,1\leqslant j\leqslant m)\) ，都有

\[
\Gamma = e _ {1} \wedge e _ {m + 1} + e _ {2} \wedge e _ {m + 2} + \dots + e _ {m} \wedge e _ {2 m}.
\]

我们称非零可分解 \(p\) -向量 \(z \in \bigwedge E\) 对 \(\Phi\) 是迷向的（相应地，全迷向的），如果向量子空间 \(V_z\) （与 \(z\) 对应；第 III 章，§ 7，第 3 小节）是迷向的（相应地，全迷向的）。

\begin{enumerate}
\def\labelenumi{\alph{enumi})}
\item
  若 z 是一个可分解 p-向量 \(\neq0\) ，2r 是 \(V_{z}\cap V_{2}^{0}\) 关于 \(V_{z}\) 的一个补的维数，证明：\(m-p+r\) 是使 \(z\wedge\Gamma^{h}\neq0\) 的整数 h 中最大的一个，其中以 \(\Gamma^{h}\) 表示 \(\Gamma\) 在外代数 \(\wedge E\) 中的 h 次幂（利用 § 4，第 2 小节的命题 2）。
\item
  证明：每个双向量 \(x \in \bigwedge E\) 都可以写成 \(\lambda \Gamma + x_{1}\) 的形式，其中 \(\lambda\) 是一个数量，而 \(x_{1}\) 是一些全迷向可分解双向量的线性组合。（化归到 \(x\) 可分解的情形，并注意若 \((e_{i})\) 是一个辛基，则 \((e_{1} + e_{2}) \wedge (e_{m+1} - e_{m+2})\) 是全迷向的。）
\item
  若 \(p \leqslant m\) ，证明：每个 \(p\) -向量 \(z \in \hat{\Lambda}\) E 都可以写成 \(z = x \wedge \Gamma + z_1\) ，其中 \(x\) 是一个 \((p - 2)\) -向量，\(z_1\) 是一些可分解且全迷向的 \(p\) -向量的线性组合。（化归到 \(z\) 可分解的情形，并对 \(p\) 作归纳。于是可化归到下述情形：设 \((e_i)\) 是一个辛基，有 \(z = e_1 \wedge e_2 \wedge \ldots \wedge e_{p-1} \wedge e_{m+p-1}\) ，并考虑双向量 \((e_{p-1} + e_{p+i}) \wedge (e_{m+p-1} - e_{m+p+i})\) （对 \(0 \leqslant i \leqslant m - p\) ）。）
\item
  对 \(1 \leqslant p \leqslant m\) ，证明：每个 \((m+p)\) -向量 \(z \in \wedge\) E 都可以写成 \(z = y \wedge \Gamma^p\) ，其中 \(y\) 是一个 \((m-p)\) -向量。（化归到 \(z\) 可分解的情形；若 \(2r\) 是 \(\mathrm{V}_z \cap \mathrm{V}_z^0\) 关于 \(\mathrm{V}_z\) 的一个补的维数，则区分两种情形：\(r = p\) 或 \(r > p\) ；在第二种情形，对 \(r\) 作归纳，其方式与 \(c)\) 中的相同。由此推出映射 \(y \to y \wedge \Gamma^p\) （从 \(\wedge^{m-p}\) E 到 \(\wedge^{m+p}\) E 中的）是双射（注意这两个空间有相同的维数）。若 y 是一个可分解的 \((m-p)\) -向量，证明 \(z = y \wedge \Gamma^{p}\) 是可分解的，并且 \(V_{z} = V_{y}^{0}\) 。
\end{enumerate}

\begin{enumerate}
\def\labelenumi{\alph{enumi})}
\setcounter{enumi}{4}
\tightlist
\item
  由 d) 推出：对 \(p \leqslant m\) ，子空间 \(\hat{E}\) 是 \(\Gamma\) 在 \(\hat{E}\) 中生成的双边理想 c 的 p 次齐次分量以及子空间 \(R_{p}\) （由适合下式的 p-向量 \(z_{1}\) 组成）的直和：
\end{enumerate}

\(z_{1} \wedge \Gamma^{m-p+1} = 0\) （对 \(z \in \bigwedge E\) ，把 \(d\) ) 应用于 \((2m - p + 2)\) -向量 \(z \wedge \Gamma^{m-p+1}\) ）。利用 \(c\) )，证明 \(R_{p}\) 由可分解的全迷向 \(p\) -向量生成，并且利用 \(d\) )，证明

映射 \(x \to x \wedge \Gamma\) （从 \(\wedge E\) 到 \(\wedge E\) 中的）是单射的，并且 \(R_p\) 的维数是 \(\binom{2m}{p} - \binom{2m}{p-2}\) 。

\begin{enumerate}
\def\labelenumi{\alph{enumi})}
\setcounter{enumi}{5}
\tightlist
\item
  证明：若 \(p \leqslant m\) ，则 p-向量 z 具有形式 \(x \wedge \Gamma\) 的一个必要充分条件是：对每个可分解的全迷向 m-向量 u，有 \(z \wedge u = 0\) 。（证明：若这个条件满足且 \(z \in R_{p}\) ，则有 \(z \wedge y = 0\) （对一切 \((2m - p)\) -向量 y）；为此，把 y 写成 \(x \wedge \Gamma^{m-p}\) 的形式，其中 x 是一个 p-向量，然后对 x 应用 c)，并注意若 \(x_{1}\) 是一个可分解的全迷向 p-向量，则 \(x_{1} \wedge \Gamma^{m-p}\) 是一个可分解的 \((2m - p)\) -向量，它可以写成 \(u_{1} \wedge \varphi_{1}\) ，其中 \(u_{1}\) 是一个可分解的全迷向 \(m\) -向量。）
\end{enumerate}

以 a 表示理想 c 在 \(\wedge\) E 中的零化子；a 是子空间 \(R_{m}\) 、\(R_{m-1} \wedge \Gamma, \ldots, R_{1} \wedge \Gamma^{m-1}\) 与 \(K\Gamma^{m}\) 的直和。证明 a 在 \(\wedge\) E 中的零化子等于 c（参见 § 2，习题 4 b））。

\begin{enumerate}
\def\labelenumi{\alph{enumi})}
\setcounter{enumi}{6}
\tightlist
\item
  对 E 的（关于 \(\Phi\) 的）每个辛自同构 u，以 \(\bar{u}\) 表示 u 到代数 \(\wedge\) E 的自同构上的典范扩张（第 III 章，§ 5，第 9 小节）。证明：\(\wedge\) E 中在所有自同构 \(\bar{u}\) 下不变的元素只有系数取自 K 的 1、\(\Gamma\) 、\(\Gamma^{2}\) 、\(\ldots\) 、\(\Gamma^{m}\) 的线性组合。（若 ( \(e_{i}\) ) 是一个辛基，写出 \(\wedge\) E 的一个元素在与 \(e_{i}\) 正交的超平面所对应的辛平延（习题 11）下不变这一条件；然后考虑辛变换 \(u_{ij}\) ，它们由 \(u_{ij}(e_{i}) = e_{j}\) ，\(u_{ij}(e_{j}) = e_{i}\) ，\(u_{ij}(e_{m+i}) = e_{m+j}\) ，\(u_{ij}(e_{m+j}) = e_{m+i}\) 以及对所有其余指标 k 有 \(u_{ij}(e_{k}) = e_{k}\) 来定义。）
\end{enumerate}

\begin{enumerate}
\def\labelenumi{\arabic{enumi})}
\setcounter{enumi}{13}
\item
  设 A 是特征 2 的交换域，E 是向量空间 \(\mathrm{A}^{2m}\) ；我们把辛群 \(\mathbf{Sp}(2m,\mathrm{A})\) 等同于辛矩阵 \(U\) 所成的群，这些矩阵满足关系 \(^t U.R.U = R\) ，其中 \(R = \left( \begin{array}{cc}0 & I_m\\ I_m & 0 \end{array} \right)\) 。令 \(D = \left( \begin{array}{cc}0 & 0\\ I_m & 0 \end{array} \right)\) ；证明：每个辛矩阵 \(U\) （使 \(I + U\) 可逆的）都可以以唯一的方式写成 \((^t D + S)^{-1}(D + S)\) 的形式，其中 \(S\) 是一个使 \(D + S\) 可逆的对称矩阵，反之亦然（参见 § 4，习题 11）。
\item
  设 A 是一个交换域，E 是 A 上维数 2m 的向量空间，Φ 是 E 上的一个非退化交错双线性型。a) 把 § 4 的习题 12 a) 与 b) 的结果推广到使 \(u^{*}u = uu^{*}\) 的 E 的自同态 u 上（\(u^{*}\) 是 u 关于 Φ 的伴随）。
\end{enumerate}

\begin{enumerate}
\def\labelenumi{\alph{enumi})}
\setcounter{enumi}{1}
\item
  我们假设 \(u^{*} = u\) 。利用 § 4 习题 12 的记号，证明：若 M 是集 \(\mathfrak{M}\) （由含于 \(G(p, p)\) 且在 \(u\) 下稳定的非迷向子空间组成）中的极小元素，则 M 或者是 \(E_{u}\) 的一个不可分解子模，或者是两个这样的同构子模的直和（按 § 4 习题 12 c) 的方式推理）。举例（取 \(p(X) = X - 1\) ）说明这两种情形都能出现。
\item
  我们假设 \(u^* = -u\) ，且 A 的特征不是 2。沿用同样的记号，设 \(p^h\) 是 \(u\) 在 M 上的限制的极小多项式，\(d\) 是 \(p\) 的次数。证明：若 \(d(h - 1)\) 是奇数，则 M 是 \(E_u\) 的一个不可分解子模；反之，若 \(d(h - 1)\) 是偶数，则 M 或者是不可分解的，或者是两个同构的不可分解子模的直和。
\item
  我们假设 \(u^{*}u = 1\) （换句话说，\(u \in \mathbf{Sp}(\Phi)\) ）。利用 c) 的记号，若 \(p(X)\) 不整除 \(X^{d(h-1)} - 1\) ，或者 \(p(X) = X - 1\) 且 \((-1)^{h} \neq -1\) ，则 M 是 \(E_{u}\) 的一个不可分解子模；否则，M 或者是不可分解的，或者是两个同构的不可分解子模的直和。
\end{enumerate}

